\documentclass[tikz,border=0pt]{standalone}
% Shared visual language for every block diagram in figures/diagrams.
% Build with XeLaTeX (see figures/diagrams/build.sh).
\usepackage{fontspec}
\setmainfont{Inter}[Numbers=Lining]
\setmonofont{DejaVu Sans Mono}[Scale=0.86]
\usepackage{xcolor}
\usetikzlibrary{arrows.meta,positioning,calc,fit,backgrounds,shapes.geometric,
  shapes.misc,matrix,decorations.pathreplacing,patterns,quotes}

% Palette (matches python/figstyle.py) --------------------------------------
\definecolor{surface}{HTML}{FCFCFB}
\definecolor{ink}{HTML}{0B0B0B}
\definecolor{ink2}{HTML}{52514E}
\definecolor{muted}{HTML}{898781}
\definecolor{hair}{HTML}{E1E0D9}
\definecolor{memline}{HTML}{2A78D6}  \definecolor{memfill}{HTML}{E7F0FC}
\definecolor{aluline}{HTML}{14875E}  \definecolor{alufill}{HTML}{E3F5EE}
\definecolor{ctlline}{HTML}{B07800}  \definecolor{ctlfill}{HTML}{FDF2D9}
\definecolor{regline}{HTML}{6250D6}  \definecolor{regfill}{HTML}{ECEAFB}
\definecolor{wire}{HTML}{3D3C39}
\definecolor{accent}{HTML}{EB6834}   \definecolor{accentfill}{HTML}{FDE8DF}
\definecolor{pos}{HTML}{E34948}      \definecolor{posfill}{HTML}{FBE1E0}
\definecolor{neg}{HTML}{2A78D6}      \definecolor{negfill}{HTML}{DCE9FA}
\definecolor{zerofill}{HTML}{F0EFEC}

\tikzset{
  every picture/.style={line cap=round, line join=round, font=\small, text=ink},
  % canvas
  canvas/.style={show background rectangle,
    background rectangle/.style={fill=surface, draw=hair, rounded corners=8pt, line width=0.6pt},
    inner frame xsep=16pt, inner frame ysep=14pt},
  ttl/.style={font=\bfseries\large, anchor=north west, inner sep=0pt, text=ink},
  sub/.style={font=\small, anchor=north west, inner sep=0pt, text=ink2},
  % blocks
  blk/.style={rounded corners=3pt, line width=0.8pt, align=center, minimum height=9mm,
    inner xsep=5pt, inner ysep=3pt},
  mem/.style={blk, draw=memline, fill=memfill},
  alu/.style={blk, draw=aluline, fill=alufill},
  ctl/.style={blk, draw=ctlline, fill=ctlfill},
  reg/.style={blk, draw=regline, fill=regfill, rounded corners=1pt},
  io/.style={blk, draw=muted, fill=white},
  ghost/.style={blk, draw=muted, dashed, fill=surface},
  % arithmetic primitives
  add/.style={circle, draw=aluline, fill=white, line width=0.8pt, minimum size=6.2mm,
    inner sep=0pt, font=\normalsize},
  shf/.style={rectangle, rounded corners=1.5pt, draw=aluline, fill=white, line width=0.7pt,
    font=\ttfamily\footnotesize, inner xsep=2.5pt, inner ysep=2pt, minimum height=5mm},
  gain/.style={isosceles triangle, isosceles triangle apex angle=60, draw=aluline, fill=white,
    line width=0.7pt, inner sep=1.2pt, font=\footnotesize},
  dn/.style={circle, draw=wire, fill=white, line width=0.7pt, minimum size=8mm, inner sep=0pt,
    font=\footnotesize},
  mux/.style={trapezium, trapezium angle=72, shape border rotate=270, draw=ctlline, fill=ctlfill,
    line width=0.8pt, minimum height=13mm, minimum width=5mm, inner sep=1pt, font=\scriptsize},
  net/.style={draw=memline, fill=memfill, rounded corners=1.5pt, font=\footnotesize, inner xsep=3pt,
    inner ysep=1.6pt, minimum width=5.5mm, minimum height=4.8mm},
  % wires
  sig/.style={-{Latex[length=2.2mm,width=1.6mm]}, draw=wire, line width=0.75pt},
  bus/.style={-{Latex[length=2.6mm,width=2mm]}, draw=wire, line width=1.5pt},
  ctrl/.style={-{Latex[length=1.8mm,width=1.4mm]}, draw=ctlline, line width=0.65pt, dash pattern=on 3pt off 2pt},
  nohead/.style={-, draw=wire, line width=0.75pt},
  dot/.style={circle, fill=wire, inner sep=0pt, minimum size=3.2pt},
  % labels
  lbl/.style={font=\footnotesize, text=ink2, fill=surface, inner sep=1.2pt},
  lblc/.style={font=\footnotesize, text=ink2, inner sep=1.2pt},
  note/.style={font=\footnotesize, text=ink2, align=left, inner sep=0pt},
  head/.style={font=\bfseries\small, text=ink, inner sep=0pt},
  code/.style={font=\ttfamily\footnotesize, text=ink},
  group/.style={draw=hair, fill=white, rounded corners=6pt, line width=0.8pt, inner sep=7pt},
  grouplabel/.style={font=\bfseries\footnotesize, text=muted, anchor=south west, inner sep=2pt},
}

% Layers: groups sit above the canvas background and below the main layer.
\pgfdeclarelayer{groups}
\pgfsetlayers{background,groups,main}

% A register: box with a clock wedge at the bottom edge.
\newcommand{\clk}[1]{\draw[regline, line width=0.6pt] ($(#1.south)+(-1.4mm,0)$) -- ($(#1.south)+(0,1.6mm)$) -- ($(#1.south)+(1.4mm,0)$);}
% Bus width tick: \bw{(point)}{width}
\newcommand{\bw}[2]{\draw[wire, line width=0.6pt] ($#1+(-0.9mm,-1.2mm)$) -- ($#1+(0.9mm,1.2mm)$); \node[font=\scriptsize, text=ink2, above=0.4mm, inner sep=0.5pt] at #1 {#2};}
\newcommand{\var}[1]{\textit{#1}}
\newcommand{\usub}[1]{\ifcase#1 ₀\or ₁\or ₂\or ₃\or ₄\or ₅\fi}
\AtBeginDocument{\hyphenpenalty=10000\exhyphenpenalty=10000}

\begin{document}
\begin{tikzpicture}[canvas]
  \node[ttl] at (0,0) {Bit-true worked example (4 × 4)};
  \node[sub] at (0,-0.62) {Integer numerators exactly as the RTL computes them: scale 8 after the row pass, 64 after the column pass. Periodic indices wrap at every border.};
  \tikzset{mx/.style={matrix of nodes, nodes={minimum width=11.5mm, minimum height=7mm, font=\small, anchor=center},
            column sep=-\pgflinewidth, row sep=-\pgflinewidth, inner sep=0pt},
           cell/.style={draw=surface, line width=1pt}}
  \matrix[mx, nodes={cell, fill=memfill}, anchor=north west] (A) at (0.2,-2.0) {
    10 & 20 & 30 & 40 \\ 50 & 60 & 70 & 80 \\ 90 & 100 & 110 & 120 \\ 130 & 140 & 150 & 160 \\};
  \node[head, anchor=south west] at ($(A.north west)+(0,0.12)$) {Input \var{x} (8-bit)};
  \matrix[mx, nodes={cell, fill=memfill}, anchor=north west] (B) at (7.0,-2.0) {
    120 & 280 & |[fill=negfill]| 0 & |[fill=negfill]| 160 \\ 440 & 600 & |[fill=negfill]| 0 & |[fill=negfill]| 160 \\
    760 & 920 & |[fill=negfill]| 0 & |[fill=negfill]| 160 \\ 1080 & 1240 & |[fill=negfill]| 0 & |[fill=negfill]| 160 \\};
  \node[head, anchor=south west] at ($(B.north west)+(0,0.12)$) {After the row pass, ×8  [L | H]};
  \matrix[mx, nodes={cell, fill=memfill}, anchor=north west] (C) at (13.8,-2.0) {
    2240 & 3520 & |[fill=negfill]| 0 & |[fill=negfill]| 1280 \\ 7360 & 8640 & |[fill=negfill]| 0 & |[fill=negfill]| 1280 \\
    |[fill=accentfill]| 0 & |[fill=accentfill]| 0 & |[fill=alufill]| 0 & |[fill=alufill]| 0 \\ |[fill=accentfill]| 5120 & |[fill=accentfill]| 5120 & |[fill=alufill]| 0 & |[fill=alufill]| 0 \\};
  \node[head, anchor=south west] at ($(C.north west)+(0,0.12)$) {Packed output, ×64  [LL LH; HL HH]};
  \draw[bus] ($(A.east)+(0.25,0)$) -- node[note, above, align=center] {rows:\\L8, H8} ($(B.west)+(-0.25,0)$);
  \draw[bus] ($(B.east)+(0.25,0)$) -- node[note, above, align=center] {columns:\\L8, H8} ($(C.west)+(-0.25,0)$);
  % highlight windows
  \draw[accent, line width=1.6pt, rounded corners=2pt] (A-1-1.north west) rectangle (A-1-2.south east);
  \draw[accent, line width=1.6pt, rounded corners=2pt, dashed] (A-1-3.north west) rectangle (A-1-4.south east);
  \draw[accent, line width=1.6pt, rounded corners=2pt] (B-1-1.north west) rectangle (B-1-1.south east);
  \draw[regline, line width=1.6pt, rounded corners=2pt] (B-1-1.north west) rectangle (B-4-1.south east);
  \draw[regline, line width=1.6pt, rounded corners=2pt] (C-1-1.north west) rectangle (C-1-1.south east);
  % equations
  \node[note, anchor=north west, text width=12.6cm] at (0.2,-5.35) {%
    \textbf{\color{accent}Row 0, pair 0.} Window \var{a}…\var{e} = \var{x}[0, −2…2] = 30, 40, 10, 20, 30 (the dashed cells are the wrapped neighbours):\\
    \texttt{L8 = −30 + 2·40 + 6·10 + 2·20 − 30 = 120}\qquad \texttt{H8 = 4·(2·20 − 10 − 30) = 0}\\
    Lifting gives the same: \texttt{dm = 80 − 30 − 10 = 40}, \texttt{dp = 40 − 10 − 30 = 0}, \texttt{L8 = 8·10 + 40 + 0 = 120}.\\[5pt]
    \textbf{Row 0, pair 1.} Window 10, 20, 30, 40, 10:\quad \texttt{L8 = 280}, \texttt{H8 = 4·(80 − 30 − 10) = 160}.};
  \node[note, anchor=north west, text width=12.2cm] at (13.2,-5.35) {%
    \textbf{\color{regline}Column 0, pair 0.} Window = rows 2, 3, 0, 1, 2 of the first packed column = 760, 1080, 120, 440, 760:\\
    \texttt{L8 = −760 + 2·1080 + 6·120 + 2·440 − 760 = 2240}\\
    \texttt{H8 = 4·(2·440 − 120 − 760) = 0}\\[5pt]
    Dividing by 64 gives the floating-point coefficients: LL = [35 55; 115 135], LH = [0 20; 0 20], HL = [0 0; 80 80], HH = 0. The output port emits 2240 first (index 0) and 0 last (index 15).};
\end{tikzpicture}
\end{document}
